%%% Кодировка и язык
\usepackage[T2A]{fontenc}
\usepackage[utf8]{inputenc}
\usepackage[english,russian]{babel}

%%% Шрифты
\usepackage{tempora}  % Times-подобный шрифт с поддержкой кириллицы
\usepackage[bigdelims,vvarbb]{newtxmath}

%%% Геометрия страницы (требования МГУ)
\usepackage[
    left=30mm,
    right=15mm,
    top=20mm,
    bottom=20mm
]{geometry}

%%% Межстрочный интервал
\usepackage{setspace}
\onehalfspacing

%%% Абзацный отступ
\usepackage{indentfirst}
\setlength{\parindent}{1.25cm}

%%% Заголовки глав и разделов
\usepackage{titlesec}
\titleformat{\chapter}[block]
    {\normalfont\Large\bfseries\centering}
    {\thechapter.~}{0pt}{}
\titleformat{\section}[block]
    {\normalfont\large\bfseries}
    {\thesection.~}{0pt}{}
\titleformat{\subsection}[block]
    {\normalfont\normalsize\bfseries}
    {\thesubsection.~}{0pt}{}

\titlespacing*{\chapter}{0pt}{-30pt}{20pt}
\titlespacing*{\section}{0pt}{12pt}{6pt}
\titlespacing*{\subsection}{0pt}{8pt}{4pt}

%%% Нумерация страниц
\usepackage{fancyhdr}
\pagestyle{fancy}
\fancyhf{}
\fancyhead[C]{\thepage}
\renewcommand{\headrulewidth}{0pt}

%%% Содержание
\usepackage{tocloft}
\renewcommand{\cfttoctitlefont}{\hfill\Large\bfseries}
\renewcommand{\cftaftertoctitle}{\hfill}
\renewcommand{\cftchapfont}{\normalfont}
\renewcommand{\cftchappagefont}{\normalfont}
\renewcommand{\cftchapleader}{\cftdotfill{\cftdotsep}}

%%% Математика
\usepackage{amsmath,amssymb,amsthm}
\usepackage{mathtools}

%%% Графика и рисунки
\usepackage{graphicx}
\graphicspath{{images/}}
\usepackage{float}
\usepackage{caption}
\usepackage{subcaption}
\captionsetup{labelsep=endash, justification=centering}

%%% Таблицы
\usepackage{booktabs}
\usepackage{longtable}
\usepackage{multirow}
\usepackage{array}

%%% Листинги кода
\usepackage{listings}
\lstset{
    basicstyle=\small\ttfamily,
    breaklines=true,
    frame=single,
    numbers=left,
    numberstyle=\tiny,
    tabsize=4,
    inputencoding=utf8,
    extendedchars=true,
    literate={а}{{\cyra}}1 {б}{{\cyrb}}1 {в}{{\cyrv}}1 {г}{{\cyrg}}1
             {д}{{\cyrd}}1 {е}{{\cyre}}1 {ж}{{\cyrzh}}1 {з}{{\cyrz}}1
             {и}{{\cyri}}1 {й}{{\cyrishrt}}1 {к}{{\cyrk}}1 {л}{{\cyrl}}1
             {м}{{\cyrm}}1 {н}{{\cyrn}}1 {о}{{\cyro}}1 {п}{{\cyrp}}1
             {р}{{\cyrr}}1 {с}{{\cyrs}}1 {т}{{\cyrt}}1 {у}{{\cyru}}1
             {ф}{{\cyrf}}1 {х}{{\cyrh}}1 {ц}{{\cyrc}}1 {ч}{{\cyrch}}1
             {ш}{{\cyrsh}}1 {щ}{{\cyrshch}}1 {э}{{\cyrerev}}1
             {ю}{{\cyryu}}1 {я}{{\cyrya}}1
}

%%% Библиография (biblatex + biber)
\usepackage[
    backend=biber,
    style=gost-numeric,
    sorting=nyt,
    language=autobib,
    autolang=other
]{biblatex}
\addbibresource{references.bib}

%%% Гиперссылки 
\usepackage[
    colorlinks=true,
    linkcolor=black,
    citecolor=black,
    urlcolor=blue
]{hyperref}

%%% ============================================================
%%% Данные о работе 
%%% ============================================================

\newcommand{\WorkType}{Магистерская диссертация}
\newcommand{\WorkTitle}{Исследование и разработка алгоритма построения рекомендаций с использованием контекстной и сессионной информации пользователя}
\newcommand{\AuthorName}{Бадекин Илья Иванович}
\newcommand{\AuthorGroup}{XXX}
\newcommand{\SupervisorName}{Ильюшин Е.А.}
\newcommand{\SupervisorTitle}{Ассистент}
\newcommand{\ReviewerName}{Фамилия И.О.}
\newcommand{\ReviewerTitle}{к.ф.-м.н., доцент}
\newcommand{\DepartmentName}{кафедра информационной безопасности}
\newcommand{\FacultyName}{факультет вычислительной математики и кибернетики}
\newcommand{\UniversityName}{Московский государственный университет имени М.В.~Ломоносова}
\newcommand{\MasterProgram}{Название магистерской программы}
\newcommand{\Year}{2026}
